% ---------------------------------------------------------------------
\section{Discussion}\label{sec:discussion}
% ---------------------------------------------------------------------

\subsection{Interpretation}\label{sec:interpretation}

% DRAFT. The interpretation of results must be written by the authors
% (SAUCIS AI policy). The prompts below mark what each paragraph needs.

\TODO{Paragraph 1: state whether weights trained at 10 and 20 chargers transfer to 200 chargers (\Fig{fig:zeroshot}, \Tab{tab:t1}) and compare the gaps}.
The main finding on zero-shot transfer is \RES{one-sentence summary}.

\TODO{Paragraph 2: relate the cost of the guarantee (\Fig{fig:cost}) to the violations of the penalty-trained policy, i.e.\ lost profit versus profit from violations}.
The main finding on the cost of the guarantee is \RES{one-sentence summary}.

\TODO{Paragraph 3: discuss how often A5 failed and the shortfall left by the fallback rule (\Fig{fig:constraints}), relating it to the closed-loop nature of A5}.
The main finding on infeasible steps is \RES{one-sentence summary}.

\TODO{Paragraph 4: compare with \Tab{tab:positioning} only on the literature configuration, never across different simulators or scenarios}.

\subsection{Practical implications for charge point operators}\label{sec:implications}

This subsection is to be written by the authors after the experiments, covering the following points.
\begin{itemize}
\item \TODO{What a CPO gains when the guarantees hold for any raw action, e.g.\ for certification or grid-operator agreements.}
\item \TODO{The cost of the guarantee in euros per EV and per year for a station of given size.}
\item \TODO{Use of the conformal variant where drivers do not declare departure times; acceptable coverage levels.}
\item \TODO{Inference time of \Tab{S-tab:t6} relative to the 15-min control step.}
\end{itemize}

\subsection{Limitations}\label{sec:limitations}

The guarantee of Theorem~\ref{thm:invariance} rests on known departure times and requested energies (A2).
Our conformal extension relaxes only the departure step, and it gives a marginal coverage statement rather than a per-EV guarantee.
Requested energies are still assumed to be declared at arrival, and strategic misreporting of deadlines \cite{ferragut2022scheduling} lies outside the model.

Aggregate feasibility (A5) can fail in closed loop, as the two-EV example of \Sec{sec:invariance} shows.
On such steps the fallback rule keeps the transformer limit and spreads the deficit in order of laxity, so the departure guarantee holds only for the EVs covered by Corollary~\ref{cor:local}.
Lemma~\ref{S-lem:deficit} bounds the remaining shortfall but does not remove it.
A layer that certifies joint feasibility over the future would need time-coupled constraints, and its projection would lose the single-sum form.

The model covers one station behind one transformer, and it ignores voltage limits, line flows and other distribution-network constraints.
Several transformers or a network model would add further coupled sum constraints to $\mathcal K(t)$, a case that Proposition~\ref{prop:projection} does not cover.

Our experiments use EV2Gym's lossless accounting, in which the billed power equals the battery-side energy change.
V2G arbitrage therefore carries no round-trip loss, so profits from discharging may be optimistic.
Efficiencies are constant per EV model instead of following the current-dependent tables of the simulator.
The theory covers lossy, piecewise-linear energy maps, but the experiments do not exercise them.

Assumption A1 requires a monotone maximum-charging map.
Lemma~\ref{S-lem:taper} proves A1 for the exponential taper of EV2Gym, provided the charging power of the EV does not exceed the charger rating.
A taper whose power drops faster than $1/(\eta^{\sm c}\Delta t)$ kW per kWh would violate A1, and certified thresholds would not repair the guarantee.
Exact delivery of sub-maximal commands (A3) also depends on the default taper multiplier of the simulator.

All results come from simulation with synthetic arrivals drawn from EV2Gym's statistics and with Dutch day-ahead prices.
Some late evening sessions of the public and private profiles do not fit into the episode horizon and are therefore not generated.
Forecast noise is synthetic, and real metering delays are not modelled.

Our EV-GNN baseline is trained per arrival profile at a single size with its own algorithm, so that comparison is not fully like-for-like.
The same holds for the MPC beyond 25 chargers, a related but not identical port.
In addition, the perfect-information LP relaxes the charging taper to the nominal power, so the reported optimality gap overestimates the true gap.
Part of any gap is also the price of not knowing the future, because the LP ignores the information limits of an online controller.
Finally, the novelty claims depend on the full texts of the closest works listed as open risks in the related-work review.
\TODO{Remove or adapt the last sentence once the novelty risks are closed}.
